\section{Prosedur dan Penyampaian Data}
\label{sec:prosedur}

Prosedur adalah sekelompok instruksi bernama yang dapat dipanggil
dari bagian lain program. Prosedur membuat program lebih mudah
dibaca, diuji, dan digunakan ulang.

Bagian ini diringkas dari \cite{gunawan2019}, Bagian 13.1,
hlm.~118--121, dan \cite{sto2001}, BAB XVI, hlm.~85--87.

\subsection{PROC, ENDP, CALL, dan RET}

Prosedur didefinisikan dengan \asm{PROC} dan diakhiri \asm{ENDP}.
Pemanggil menggunakan \asm{CALL}, sedangkan prosedur mengakhiri
dirinya dengan \asm{RET}.

\begin{lstlisting}[language=MASM, numbers=none]
HitungLuas PROC
    ; isi prosedur
    RET
HitungLuas ENDP

; program utama
CALL HitungLuas
\end{lstlisting}

\asm{CALL} menyimpan alamat instruksi sesudah \asm{CALL} ke stack,
lalu memindahkan kendali ke prosedur. \asm{RET} mengambil alamat itu
dari stack dan memuatnya kembali ke \reg{IP}.

\begin{table}[H]
    \centering
    \caption{Pemanggilan dekat dan jauh pada 8086.}
    \label{tab:call-near-far}
    \small
    \begin{tabular}{@{}p{3.2cm}p{5.0cm}p{5.0cm}@{}}
        \toprule
        \textbf{Aspek} & \textbf{\asm{CALL} dekat} &
        \textbf{\asm{CALL} jauh} \\
        \midrule
        Alamat yang disimpan & \reg{IP} (2 byte) &
        \reg{CS} dan \reg{IP} (4 byte) \\
        \addlinespace
        Instruksi kembali & \asm{RET} & \asm{RETF} \\
        \addlinespace
        Pemakaian & Prosedur pada segmen kode yang sama &
        Prosedur pada segmen kode lain \\
        \bottomrule
    \end{tabular}
\end{table}

Seluruh contoh pada bab ini memakai pemanggilan dekat karena
prosedur dan pemanggil berada pada satu segmen kode.

\subsection{Stack Frame dan Register BP}

Parameter yang dikirim melalui stack dapat diakses dengan register
\reg{BP}. Sebelum mengakses parameter, prosedur menyimpan \reg{BP}
lama dan menyalin \reg{SP} ke \reg{BP}:

\begin{lstlisting}[language=MASM, numbers=none]
PUSH BP             ; simpan BP pemanggil
MOV BP, SP          ; BP menunjuk puncak stack frame
\end{lstlisting}

Setelah itu, isi stack frame dapat diakses dengan mode
\textit{base relative}:

\begin{itemize}
    \item \code{[BP+0]}: nilai \reg{BP} lama;
    \item \code{[BP+2]}: alamat kembali;
    \item \code{[BP+4]}: parameter pertama;
    \item \code{[BP+6]}: parameter kedua.
\end{itemize}

Susunan stack frame ditunjukkan pada Gambar~\ref{fig:stack-frame}.

\begin{figure}[H]
    \centering
    \begin{tikzpicture}[
        sel/.style={
            draw=obeBlue, line width=0.7pt, minimum width=4.6cm,
            minimum height=0.65cm, align=center, font=\small,
            inner sep=2pt
        },
        param/.style={sel, fill=obeAmberLight, draw=obeAmber},
        kendali/.style={sel, fill=obeTealLight, draw=obeTeal},
        lama/.style={sel, fill=obeGrayLight, draw=obeGray},
        ket/.style={font=\scriptsize, text=obeGray, anchor=west}
    ]
        \node[param] (p2) at (0,0) {parameter kedua \quad [BP+6]};
        \node[param] (p1) [below=0pt of p2] {parameter pertama \quad [BP+4]};
        \node[kendali] (ret) [below=0pt of p1]
            {alamat kembali \quad [BP+2]};
        \node[lama] (bp) [below=0pt of ret] {BP lama \quad [BP+0]};
        \node[sel, fill=obeBlueLight] (lokal) [below=0pt of bp]
            {variabel lokal / ruang bebas};

        \node[ket] at (2.45,0) {alamat tinggi};
        \node[ket] at (2.45,-3.25) {alamat rendah};
        \node[font=\scriptsize, text=obeMaroon, anchor=east]
            at (-2.5,-1.95) {BP dan SP};
        \draw[-{Stealth[length=2.2mm]}, line width=0.9pt, obeMaroon]
            (-2.5,-1.95) -- (-2.35,-1.95);
    \end{tikzpicture}
    \caption{Stack frame prosedur: parameter, alamat kembali, BP lama,
    dan ruang lokal.}
    \label{fig:stack-frame}
\end{figure}

\subsection{Penyampaian Data}

Data dapat disampaikan ke prosedur melalui tiga cara.

\begin{table}[H]
    \centering
    \caption{Cara penyampaian data ke prosedur.}
    \label{tab:parameter}
    \small
    \begin{tabular}{@{}p{2.6cm}p{4.6cm}p{5.8cm}@{}}
        \toprule
        \textbf{Cara} & \textbf{Contoh} & \textbf{Keterangan} \\
        \midrule
        Register & \lstasm|MOV AX, lebar| lalu \lstasm|CALL Hitung| &
        Cepat; jumlah parameter terbatas \\
        \addlinespace
        \eng{Stack} & \lstasm|PUSH lebar| lalu \lstasm|CALL Hitung| &
        Cocok untuk banyak parameter; prosedur mengakses
        \code{[BP+4]} \\
        \addlinespace
        Variabel & \lstasm|MOV lebar, 5| lalu \lstasm|CALL Hitung| &
        Sederhana; prosedur bergantung pada variabel global \\
        \bottomrule
    \end{tabular}
\end{table}

Bila parameter dikirim melalui stack, pemanggil mendorong parameter
sebelum \asm{CALL}. Prosedur dapat membersihkan parameter dengan
\asm{RET} \code{n}, dengan $n$ adalah jumlah byte parameter.

\begin{lstlisting}[language=MASM, numbers=none]
PUSH tinggi         ; parameter kedua
PUSH lebar          ; parameter pertama
CALL HitungLuas
; prosedur memakai RET 4 untuk membuang dua parameter
\end{lstlisting}

\begin{peringatanbox}
\vspace{1.1em}
Prosedur sebaiknya menyimpan register yang dipakainya dengan
\asm{PUSH} di awal dan memulihkannya dengan \asm{POP} di akhir.
Tanpa itu, pemanggil dapat kehilangan nilai register yang masih
dibutuhkan. Pastikan pula jumlah \asm{PUSH} dan \asm{POP} seimbang.
\end{peringatanbox}
